problem:  
Let \( \pi : (M^{4n}, g) \to (B^{2n}, g_B) \) be a **Riemannian submersion** between closed, simply connected Riemannian manifolds.  
Assume that both \(M\) and \(B\) have **positive sectional curvature**, and that all fibers \(F_b = \pi^{-1}(b)\) are **totally geodesic** and compact of dimension \(2n\) (thus \( \dim M = 4n\)).  

Denote by \(R^M\), \(R^B\), and \(R^F\) their respective curvature tensors, and by \(A\) the **O’Neill integrability tensor**. The curvature of \(M\) can be represented as
\[
R^M = \pi^*R^B + R^F + \mathcal{R}(A),
\]
where \(\mathcal{R}(A)\) denotes the part produced by \(A\) via O’Neill’s formulas.

Let \(p_1(M)\) denote the first Pontryagin form on \(M\), constructed via Chern–Weil theory from \(R^M\), and let \(L(p_1(M), \ldots)\) be the corresponding **Hirzebruch \(L\)-form**.

**Question:**  
If both the total space \(M\) and base \(B\) have positive sectional curvature, must the **Hirzebruch signature**
\[
\sigma(M) = \langle L(p_1(M), \ldots), [M]\rangle
\]
be strictly positive?  
Equivalently, can there exist a positively curved Riemannian submersion with totally geodesic fibers in which the mixed curvature terms from \(A\) cause the Pontryagin form integrals to cancel to zero or to yield a nonpositive signature?

Why is it a "good" problem:  
This problem avoids the topological triviality of Euler characteristic multiplicativity by focusing on the **signature**, a subtle topological invariant tied to the curvature tensor through the **Chern–Weil construction** but with no known multiplicative behavior under submersion. The question probes whether the positivity of sectional curvature, globally enforced on both total and base spaces, constrains the sign of a purely topological quantity built from **traces of products of curvature forms**.  

Unlike Euler characteristic, the signature depends on the *algebraic sign* structure within the curvature operator, precisely where O’Neill’s \(A\)-tensor contributes mixed terms not captured by base and fiber geometry alone. Whether such terms can interfere destructively with Chern–Weil densities is completely unknown.  

Thus, the problem bridges differential geometry, algebraic topology, and global analysis. It asks for a new kind of **curvature-sign rigidity**: does global positive curvature enforce topological “positivity” of oriented invariants? A resolution either way would represent a major advance in understanding the interplay between curvature conditions and characteristic class integrals—a frontier area not addressed by existing techniques.